\section{Hallucination Definition and Taxonomy}

Hallucination is one of the central challenges facing MLLMs.
While different studies adopt varying perspectives on definition and categorization,
a broad consensus has emerged.
This section reviews three representative surveys
and then presents the taxonomy adopted in this work
\citep{baiHallucinationMultimodalLarge2025,
       chenSurveyMultimodalHallucination2026,
       liuSurveyHallucinationLarge2024}.

\citet{baiHallucinationMultimodalLarge2025} define hallucination in MLLMs
as generated textual responses that are inconsistent
with the corresponding visual content,
emphasizing cross-modal inconsistency.
They also trace the concept back to LLMs,
distinguishing factual hallucinations (deviations from real-world facts)
from faithfulness hallucinations
(deviations from input context or self-consistency).
Their object-centric taxonomy comprises three levels:
category (fabricating or misidentifying objects),
attribute (correct category but wrong properties such as color, shape, or count),
and relation (correct objects and attributes
but incorrect spatial or interactive relations).

\citet{chenSurveyMultimodalHallucination2026} propose a broader unified definition,
treating hallucination as any inconsistency
between generated content and either the input conditions
or established world knowledge.
Their taxonomy extends the faithfulness-factuality dichotomy:
faithfulness hallucinations are stratified by visual granularity
into object-level, attribute-level, and scene-level errors,
forming a progression from shallow perception to deep understanding;
factuality hallucinations are divided into commonsense errors
and domain-specific errors.

\citet{liuSurveyHallucinationLarge2024} restrict hallucination
to the visual domain,
defining it as inconsistency between factual visual content in the image
and the corresponding generated text.
They adopt a dual-perspective taxonomy:
a cognitive perspective distinguishing judgment hallucinations
(yes/no contradictions with visual facts)
from descriptive hallucinations
(unfaithful free-form generation),
and a semantic perspective distinguishing object, attribute,
and relation hallucinations.

Synthesizing these frameworks and considering that our experiments
primarily target visual question answering and image-text reasoning tasks,
we adopt the following definition and three-category taxonomy.

\begin{quote}
\textit{We define MLLM hallucination as a model response
that is inconsistent with the input visual content,
contradicts verifiable real-world knowledge,
or contains errors in the reasoning process.}
\end{quote}

\noindent The three categories are detailed below.

\subsection{Faithfulness Hallucination (Visual Inconsistency)}

Faithfulness hallucination refers to generated text
that is inconsistent with the visual content of the input image.
Depending on the level of visual granularity,
this category is further divided into three levels.

At the object level,
the model claims the presence of objects that do not exist in the image,
omits objects that are actually present,
or misidentifies object categories.
The POPE benchmark systematically detects this type of hallucination
through binary probing questions
(\Cref{fig:faith_object}).

At the attribute level,
the model correctly identifies the object category
but incorrectly describes its properties
such as color, shape, material, or pose
(\Cref{fig:faith_attr}).

At the scene level,
spatial relations, interactive events,
or other scene-level descriptions are inconsistent with the image,
even though all individual objects are correctly identified
(\Cref{fig:faith_scene}).

\begin{figure}[htbp]
  \centering
  \begin{subfigure}[t]{0.5\textwidth}
    \centering
    \includegraphics[width=\linewidth]{faith_object.pdf}
    \caption{Object level.}
    \label{fig:faith_object}
  \end{subfigure}\hfill
  \begin{subfigure}[t]{0.45\textwidth}
    \centering
    \includegraphics[width=\linewidth]{faith_attr.pdf}
    \caption{Attribute level.}
    \label{fig:faith_attr}
  \end{subfigure}\hfill
  \begin{subfigure}[t]{0.5\textwidth}
    \centering
    \includegraphics[width=\linewidth]{faith_scene.pdf}
    \caption{Scene level.}
    \label{fig:faith_scene}
  \end{subfigure}
  \caption{Faithfulness hallucination examples at three levels of visual granularity.}
  \label{fig:faith_examples}
\end{figure}

The central detection question for this category is:
is the output consistent with the image?

\subsection{Factuality Hallucination (Factual Inconsistency)}

Factuality hallucination refers to generated outputs
that contradict established, verifiable real-world knowledge.

At the commonsense level,
the error involves everyday knowledge,
where the model's visual perception is roughly correct
but the output contradicts well-known facts
(\Cref{fig:fact_common}).

At the domain-specific level,
the error involves specialized professional knowledge.
Such hallucinations are particularly common
in the VQA-RAD dataset
(\Cref{fig:fact_domain}).

\begin{figure}[htbp]
  \centering
  \begin{subfigure}[t]{0.48\textwidth}
    \centering
    \includegraphics[width=\linewidth]{fact_common.pdf}
    \caption{Commonsense level.}
    \label{fig:fact_common}
  \end{subfigure}\hfill
  \begin{subfigure}[t]{0.48\textwidth}
    \centering
    \includegraphics[width=\linewidth]{fact_attr.pdf}
    \caption{Domain-specific level.}
    \label{fig:fact_domain}
  \end{subfigure}
  \caption{Factuality hallucination examples
    at the commonsense and domain-specific levels.}
  \label{fig:fact_examples}
\end{figure}

The central detection question for this category is:
is the output consistent with world knowledge?

\subsection{Logical Hallucination (Reasoning Error)}

Logical hallucination refers to errors in the reasoning process,
where the model's conclusion does not logically follow
from the visual evidence or from its own established premises.

In deductive reasoning errors,
the model correctly reads the visual information
but arrives at a contradictory conclusion.
Such numerical reasoning errors
are the primary source of logical hallucinations in MathVista
(\Cref{fig:logic_deductive}).

In causal reasoning errors,
the model infers causal relationships
that contradict the visual evidence
(\Cref{fig:logic_causal}).

\begin{figure}[htbp]
  \centering
  \begin{subfigure}[t]{0.48\textwidth}
    \centering
    \includegraphics[width=\linewidth]{logic_deductive.pdf}
    \caption{Deductive reasoning.}
    \label{fig:logic_deductive}
  \end{subfigure}\hfill
  \begin{subfigure}[t]{0.48\textwidth}
    \centering
    \includegraphics[width=\linewidth]{logic_causal.pdf}
    \caption{Causal reasoning.}
    \label{fig:logic_causal}
  \end{subfigure}
  \caption{Logical hallucination examples
    in deductive and causal reasoning.}
  \label{fig:logic_examples}
\end{figure}

The central detection question for this category is:
is the reasoning valid?
